\section{Implementation}
\label{sec:implementation}

This section describes how each feature is implemented, from the database to the interface.

\subsection{F1 -- Accounts and administration}

\texttt{UserService} creates the accounts (role \texttt{user}, username and e-mail unique) and
checks the passwords at login. A successful login returns the two tokens and is recorded in the
\texttt{connection} table. The frontend keeps the tokens in the Streamlit session; when the access
token expires, the HTTP client asks for a new one with the refresh token and sends the request
again, so the user is not logged out after 30 minutes.

The administrators have a dedicated page. It lists the accounts, promotes a user to the
\texttt{admin} role, deletes an account (with its favourites and alerts, through
\texttt{ON DELETE CASCADE}) and shows the connection history of each user. An administrator
cannot delete their own account, so that at least one administrator always remains. The menu of
an administrator also shows the date of the last import, a button to start an update and a
button to reset the database.

\subsection{F2 -- Dashboard}

\texttt{DashboardDao} runs four queries on a period chosen by the user (by default the next seven
days): the approaches of the period, the number of distinct NEOs and of hazardous ones, the
closest approach, and the next approach of each favourite of the user. \texttt{DashboardService}
checks the period and computes the share of hazardous objects.

\subsection{F3 -- Search and details}

\texttt{NeoDao.find} builds the query from the criteria that are given: each criterion adds one
condition with a named parameter, and the values are always passed to \texttt{psycopg2} as
parameters, never written into the SQL text, which prevents SQL injection. The \texttt{approach}
table is only joined when a date criterion is used. The sort column cannot be a SQL parameter, so
only the columns of a whitelist are accepted. The detailed view lists the approaches with their
distance in kilometres and in astronomical units; the results and the approaches can be exported
as CSV, and a NEO can be added to the favourites from this view.

\subsection{F4 -- Declaring a NEO}

A logged-in user gives a name and, optionally, the absolute magnitude, the estimated diameters
and the hazard flag. The values are checked twice: by the Pydantic schema (positive diameters,
minimum not above maximum) and by \texttt{NeoService}, so that the rule also holds for any other
caller of the service. The NEO is stored with origin \texttt{USER} and no NASA identifier.

\subsection{F5 -- Favourites}

\texttt{FavouriteService} checks that the NEO exists before adding it, and an object cannot be
added twice (primary key \texttt{(user\_id, neo\_id)}). For each favourite, the page shows the
list of its approaches, exportable as CSV, and a graph of its distance to the Earth over time.
The graph is drawn by the backend with Matplotlib and sent as a PNG image encoded in base64; the
page displays it and offers it for download.

\subsection{F6 -- Alerts and notifications}

An alert combines up to three criteria: a minimum diameter, a maximum distance and ``hazardous
objects only''; at least one is required, and two alerts of the same user cannot have exactly the
same criteria. A single SQL query finds the future approaches that match the active alerts of a
user.

The notifications are evaluated when the user logs in, and each approach is notified only once:
when an alert has matched, its \texttt{last\_eval\_date} is set to the current time of the
database, and the query only keeps the approaches recorded after that date. Both dates come from
the database clock, so the comparison does not depend on the time zone of the server. The
notifications of the session stay visible on the alerts page.

\subsection{FO3 -- Statistics}

\texttt{StatisticsDao} computes, with \texttt{GROUP BY} and aggregate functions, the number of
approaches per month, the number of NEOs per size class and a summary: number of NEOs, share of
hazardous ones, mean and median distance (\texttt{PERCENTILE\_CONT}). The size classes follow the
thresholds used by NASA (25~m, 140~m, 1~km); NEOs declared without a diameter are counted in an
``unknown'' class. Each statistic is displayed with a short explanation of why it matters, for
instance why the median describes a typical approach better than the mean.

\subsection{FO7 -- Automatic update}

\texttt{DailyImportScheduler} is started with the API. It starts an import at once, then every
24~hours, in the name of the first administrator. Each run imports the days missing at the end of
the one-year window, and re-imports the next 30 days once a day to get the objects discovered
since the last import (about five calls). The administrator button calls the same code, and a
lock guarantees that only one import runs at a time; the database cannot be reset while an import
is running. The automatic import can be turned off with \texttt{AUTO\_IMPORT=false}.

\subsection{FO8 -- Picture of the day}

\texttt{ApodService} calls the NASA APOD API for the chosen day. The page shows the image and its
description and offers both for download; when NASA publishes a video instead of an image, a link
to it is shown instead. If NASA is unavailable, the API answers with a clear 503 error instead of
failing.

\subsection{API endpoints}

\begin{longtable}{l p{5.6cm} p{6.4cm}}
\caption{Endpoints of the API (full documentation: \texttt{/docs})} \\
\toprule
\textbf{Method} & \textbf{Path} & \textbf{Description} \\
\midrule
\endfirsthead
\toprule
\textbf{Method} & \textbf{Path} & \textbf{Description} \\
\midrule
\endhead
POST & \texttt{/login/} & Log in, returns the tokens \\
POST & \texttt{/login/register} & Create an account \\
POST & \texttt{/login/refresh} & New tokens from a refresh token \\
GET & \texttt{/login/me} & Profile of the logged-in user \\
GET & \texttt{/neo} & List the NEOs (pagination) \\
POST & \texttt{/neo/search} & Search with criteria \\
GET & \texttt{/neo/\{neo\_id\}} & Details and approaches of a NEO \\
POST & \texttt{/neo} & Declare a NEO (origin \texttt{USER}) \\
GET & \texttt{/dashboard/} & Dashboard of a period \\
GET & \texttt{/statistics/} & Statistics \\
GET & \texttt{/apod/} & Picture of the day \\
POST, DELETE & \texttt{/favourites/\dots} & Add and remove favourites \\
GET & \texttt{/favourites/user/\{id\}/\dots} & Favourites, their approaches, their graph \\
POST, GET, PUT, DELETE & \texttt{/alerts/\dots} & Create, list, switch on/off, delete alerts \\
GET & \texttt{/alerts/user/\{id\}/notifications} & New notifications \\
POST & \texttt{/import/update} & Start the NASA update (admin) \\
GET & \texttt{/import/status}, \texttt{/import/last} & Running import, end of the last import \\
POST & \texttt{/reset\_database} & Re-create the database (admin) \\
GET & \texttt{/users/} & List the accounts (admin) \\
PUT & \texttt{/users/\{id\}/admin} & Promote to administrator (admin) \\
DELETE & \texttt{/users/\{id\}} & Delete an account (admin) \\
GET & \texttt{/users/\{id\}/connections} & Connection history (admin) \\
\bottomrule
\end{longtable}
